\chapter{Sprint 4 : Gestion des Endpoints et des Modèles}

\section{Introduction}
Ce sprint se focalise sur le développement de l'API RESTful et la définition précise des modèles de données Mongoose pour assurer une persistence cohérente et performante.

\section{Backlog du Sprint 4}
\begin{itemize}
    \item \textbf{US-10} : Développement des contrôleurs pour les Cours et les Quiz.
    \item \textbf{US-11} : Définition des schémas de données complexes (Attendance, Submissions).
    \item \textbf{US-12} : Mise en place des routes API sécurisées.
\end{itemize}

\section{Analyse}
\subsection{Diagramme de cas d'utilisation}
L'Enseignant crée des examens et des quiz. L'Étudiant soumet ses travaux. L'Administrateur génère des PV de notes.

\subsection{Description textuelle des cas d'utilisation}
L'enseignant définit les paramètres d'un quiz (temps, questions). L'étudiant répond aux questions. Le système enregistre les réponses dans la collection \texttt{StudentQuizAttempt} et calcule le score automatiquement.

\section{Conception}
\subsection{Diagrammes de séquence}
Le diagramme de séquence pour la soumission d'un examen illustre le flux depuis le client jusqu'à la validation et l'enregistrement des notes en base de données.

\subsection{Diagramme de classes}
Détail des modèles \texttt{Quiz}, \texttt{Question}, \texttt{Exam}, \texttt{Submission} et \texttt{Grade}.

\subsection{Représentation graphique de la base de données}
Schéma complet de la base de données MongoDB illustrant les relations entre les collections académiques et les collections d'évaluation.
\begin{figure}[H]
    \centering
    \fbox{Placeholder: Schéma ER / MongoDB}
    \caption{Architecture des données d'eduGenius}
\end{figure}

\section{Réalisation}
\subsection{Interface de Création de Quiz}
Formulaire dynamique permettant d'ajouter des questions, des options et de définir la bonne réponse.
\begin{figure}[H]
    \centering
    \fbox{Placeholder: Capture d'écran - Quiz Creator}
    \caption{Interface de création d'évaluations}
\end{figure}

\subsection{Implémentation des Contrôleurs REST}
Développement de la logique métier dans \texttt{QuizController.js} et \texttt{ExamController.js}.

\section{Tests}
\subsection{Tests unitaires}
Validation de la logique de calcul des scores des quiz.
\subsection{Tests d'intégration}
Test des endpoints de soumission de travaux avec gestion des fichiers joints.
\subsection{Tests fonctionnels}
Vérification de la génération correcte des moyennes dans le module PV.

\section{Conclusion}
L'API étant complète et robuste, le dernier sprint sera consacré à l'interface utilisateur finale et à l'intégration poussée de l'IA.
